\section{Авторизация}\label{ACME.Authz}

В процессе авторизации сервер проверяет, что
\begin{itemize}
\item 
клиент владеет личным ключом учетной записи;
\item 
клиент контролирует целевой идентификатор. 
\end{itemize}

Авторизацию можно повторить, чтобы связать с учетной записью несколько 
идентификаторов (например, для выпуска сертификата с несколькими 
идентификаторами в расширении \token{SubjectAltName}{ENTITIES.SAN}) или связать 
с идентификатором несколько учетных записей (например, для выпуска нескольких 
сертификатов для одного идентификатора).  

Объекты авторизации создаются сервером в ответ на 
\pdfexample{запрос}{acme/authz1.txt} на узел \code{newOrder} или 
\code{newAuthz}, отправленный владельцем учетной записи. URL-адреса 
объектов отправляются клиенту в ответе на запрос. Объект авторизации неявно 
связан с ключом учетной записи, который используется для подписи запроса.

При запросе на узел \code{newOrder} клиент получает от сервера заявку на выпуск 
сертификата и отправляет \pdfexample{запросы}{acme/authz2.txt} POST-as-GET на 
указанные в заявке URL-адреса, получая в \pdfexample{ответ}{acme/authz3.txt} 
объекты авторизации. 
%
При запросе на узел \code{newAuthz} клиент заранее загружает объекты 
авторизации.

\if0
\begin{codeblock}
POST /acme/authz/PAniVnsZcis HTTP/1.1
Host: example.com
Content-Type: application/jose+json

{
  "protected": base64url({
    "alg": "BIGNS128",
    "kid": "https://example.com/acme/acct/evOfKhNU60wg",
    "nonce": "uQpSjlRb4vQVCjVYAyyUWg",
    "url": "https://example.com/acme/authz/PAniVnsZcis"
  }),
  "payload": "",
  "signature": "nuSDISbWG8mMgE7H...QyVUL68yzf3Zawps"
}
\end{codeblock}
    
\begin{codeblock}
HTTP/1.1 200 OK
Content-Type: application/json
Link: <https://example.com/acme/directory>;rel="index"

{
  "status": "pending",
  "expires": "2026-01-02T14:09:30Z",

  "identifier": {
    "type": "dns",
    "value": "www.example.org"
  },

  "challenges": [
    {
      "type": "http-01",
      "url": "https://example.com/acme/chall/prV_B7yEyA4",
      "token": "DGyRejmCefe7v4NfDGDKfA",
      "status": "pending"
    },
    {
      "type": "dns-01",
      "url": "https://example.com/acme/chall/Rg5dV14Gh1Q",
      "token": "DGyRejmCefe7v4NfDGDKfA",
      "status": "pending"
    }
  ]
}
\end{codeblock}
\fi

% use: https://errata.rfc-editor.org/eid8883/

Чтобы завершить авторизацию, подтвердив контроль над идентификатором, клиенту
необходимо отработать вызов авторизации и сообщить серверу, что он готов к
проверке отработки.

Клиент сообщает серверу о готовности, отправляя 
\pdfexample{POST-запрос}{acme/authz4.txt} на URL-адрес вызова авторизации 
(а не на URL-адрес объекта авторизации). Тело запроса представляет собой объект 
JWS, в поле \sstr{payload} которого серверу передаются сведения, необходимые 
для проверки отработки вызова авторизации и обновления статуса объекта 
авторизации. Для всех типов вызовов авторизации, определяемых в настоящем 
стандарте, поле \sstr{payload} представляет собой пустой объект JSON 
(\sstr|{}|).

% use: https://errata.rfc-editor.org/eid5729/

\if0
Например, если бы клиент ответил на вызов авторизации \sstr{http-01} в примере 
выше, он бы отправил следующий запрос:  

\begin{codeblock}
POST /acme/chall/prV_B7yEyA4 HTTP/1.1
Host: example.com
Content-Type: application/jose+json

{
  "protected": base64url({
    "alg": "BIGNS128",
    "kid": "https://example.com/acme/acct/evOfKhNU60wg",
    "nonce": "Q_s3MWoqT05TrdkM2MTDcw",
    "url": "https://example.com/acme/chall/prV_B7yEyA4"
  }),
  "payload": base64url({}),
  "signature": "9cbg5JO1Gf5YLjjz...SpkUfcdPai9uVYYQ"
}
\end{codeblock}
\fi

Сервер обновляет объект авторизации, уточняя статус вызова авторизации по 
результатам его отработки клиентом. 
%
% skip: The server MUST ignore any fields in the response object
%  that are not specified as response fields for this type of challenge.
%  Note that the challenges in this document do not define any response
%  fields, but future specifications might define them.  
%
При запросе объекта авторизации сервер возвращает ответ c кодом 
статуса 200 и обновленной информацией о вызове авторизации в поле 
\sstr{challenges}.

Если ответ клиента некорректен или не дает серверу всю информацию,
необходимую для проверки отработки вызова, сервер должен возвратить ошибку.  
При получении ошибки клиенту следует отменить все действия, предпринятые для 
отработки вызова, например, удалить файлы, размещенные на веб-сервере. 

% info: В настоящем стандарте полезная часть ответа -- пустой объект.
% В будущем ответ может быть более информативным.

Сервер финализирует авторизацию, если успешно завершена проверка отработки 
одного из вызовов авторизации или если все проверки завершены с отрицательным 
результатом.
%
% use: https://errata.rfc-editor.org/eid6317/
%
Объекту авторизации присваивается статус \sstr{valid} или \sstr{invalid} 
в зависимости от того, считает ли сервер авторизацию успешной. Если присвоен 
статус \sstr{valid}, то сервер должен добавить в объект авторизации  
поле \sstr{expires}. При завершении авторизации сервер может удалять 
неотработанные вызовы авторизации и может изменять поле \sstr{expires}. 
Сервер не должен удалять вызовы со статусом \sstr{invalid}. 
 
Обычно процесс проверки отработки вызова занимает некоторое время, поэтому 
клиенту необходимо запросить объект авторизации, чтобы узнать, завершена 
ли проверка. 
%
Клиент отправляет \pdfexample{запрос}{acme/authz5.txt} POST-as-GET на адрес 
авторизации, а сервер отвечает текущим объектом авторизации. До тех пор, пока 
проверка не завершена, в \pdfexample{ответах}{acme/authz6.txt} клиенту сервер 
должен указывать код статуса~200. Сервер может включать в ответы заголовок 
\code{Retry-After}, чтобы подсказать клиенту интервал опроса. 
%
Если клиент имеет возможность определить, когда сервер проверил обработку 
(например, по HTTP- или DNS-запросу от сервера), то клиенту не следует 
запрашивать объект авторизации до момента проверки.

\if0
\begin{codeblock}
POST /acme/authz/PAniVnsZcis HTTP/1.1
Host: example.com
Content-Type: application/jose+json

{
  "protected": base64url({
    "alg": "BIGNS128",
    "kid": "https://example.com/acme/acct/evOfKhNU60wg",
    "nonce": "uQpSjlRb4vQVCjVYAyyUWg",
    "url": "https://example.com/acme/authz/PAniVnsZcis"
  }),
  "payload": "",
  "signature": "nuSDISbWG8mMgE7H...QyVUL68yzf3Zawps"
}
\end{codeblock}

\begin{codeblock}
HTTP/1.1 200 OK
Content-Type: application/json
Link: <https://example.com/acme/directory>;rel="index"

{
  "status": "valid",
  "expires": "2028-09-09T14:09:01.13Z",

  "identifier": {
    "type": "dns",
    "value": "www.example.org"
  },

  "challenges": [
    {
      "type": "http-01",
      "url": "https://example.com/acme/chall/prV_B7yEyA4",
      "status": "valid",
      "validated": "2024-12-01T12:05:13.72Z",
      "token": "IlirfxKKXAsHtmzK29Pj8A"
    }
  ]
}
\end{codeblock}
\fi

Клиент может отменить авторизацию, отправив 
\pdfexample{POST-запросы}{acme/authz7.txt} с подписанным фиксированным объектом 
\code|{"status":"deactivated"}| на URL-адреса объектов авторизации. 

\if0
\begin{codeblock}
POST /acme/authz/PAniVnsZcis HTTP/1.1
Host: example.com
Content-Type: application/jose+json

{
  "protected": base64url({
    "alg": "BIGNS128",
    "kid": "https://example.com/acme/acct/evOfKhNU60wg",
    "nonce": "xWCM9lGbIyCgue8di6ueWQ",
    "url": "https://example.com/acme/authz/PAniVnsZcis"
  }),
  "payload": base64url({
    "status": "deactivated"
  }),
  "signature": "srX9Ji7Le9bjszhu...WTFdtujObzMtZcx4"
}
\end{codeblock}
\fi

Сервер должен проверить, что запрос подписан ключом той учетной записи, 
которая касается авторизации. Если сервер согласен с деактивацией, 
то он должен дать ответ с кодом статуса 200 и обновленным объектом авторизации.

После деактивации объект авторизации не может использоваться в качестве 
основания для выпуска сертификата.
